% Opción múltiple (0.75 pts c/u). La respuesta correcta va en negritas y en rojo.

Marque una sola respuesta.

\begin{enumerate}[label=\textbf{\arabic*.},leftmargin=*,itemsep=8pt]
	
	\item El error local de un método de un paso es:
	\begin{enumerate}[label=\alph*)]
		\item el error acumulado en todos los pasos hasta $T$
		\item {\bfseries\boldmath\color{red} el error cometido en un paso suponiendo exacta la solución en el instante anterior}
		\item el residuo al sustituir la solución numérica en la ecuación
		\item el error de redondeo de la máquina
	\end{enumerate}
	
	\item Si un método de un paso estable tiene error local $\Order{h^{p+1}}$, su error global en $T$ fijo es:
	\begin{enumerate}[label=\alph*)]
		\item $\Order{h^{p+1}}$ 
		\item $\Order{h^{p+2}}$ 
		\item {\bfseries\boldmath\color{red} $\Order{h^{p}}$} 
		\item $\Order{h^{p-1}}$
	\end{enumerate}
	
	\item Para $N' = -\lambda N$, el método de Crank--Nicolson tiene error local y global, respectivamente:
	\begin{enumerate}[label=\alph*)]
		\item $\Order{h^2}$ y $\Order{h}$ 
		\item $\Order{h^2}$ y $\Order{h^2}$
		\item $\Order{h^4}$ y $\Order{h^3}$ 
		\item {\bfseries\boldmath\color{red} $\Order{h^3}$ y $\Order{h^2}$}
	\end{enumerate}
	
	\item Euler explícito para $N' = -\lambda N$ ($\lambda > 0$) es absolutamente estable si:
	\begin{enumerate}[label=\alph*)]
		\item  $0 < \lambda h < 2$ {\bfseries\boldmath\color{red}
		\item $0 < \lambda h < 1$} 
		\item $0 < \lambda h < 4$ 
		\item para todo $h > 0$
	\end{enumerate}
	
	\item La aproximación $\dfrac{f(x+h) - 2f(x) + f(x-h)}{h^2}$ de $f''(x)$ tiene error de truncamiento:
	\begin{enumerate}[label=\alph*)]
		\item $\Order{h}$ \item $\Order{h^3}$ \item $\Order{h^4}$ \item {\bfseries\boldmath\color{red} $\Order{h^2}$}
	\end{enumerate}
	
	\item El radio espectral de una matriz es:
	\begin{enumerate}[label=\alph*)]
		\item la suma de sus valores propios 
		\item {\bfseries\boldmath\color{red} el máximo módulo de sus valores propios}
		\item su determinante 
		\item su norma de Frobenius
	\end{enumerate}
	
	\item La rapidez de convergencia del método de las potencias está gobernada por:
	\begin{enumerate}[label=\alph*)]
		\item $|\lambda_1|/|\lambda_2|$ 
		\item el número de condición de $\mat{A}$
		\item {\bfseries\boldmath\color{red} $|\lambda_2|/|\lambda_1|$} 
		\item la traza de $\mat{A}$
	\end{enumerate}
	
	\item El método de Wielandt acelera la convergencia porque:
	\begin{enumerate}[label=\alph*)]
		\item {\bfseries\boldmath\color{red} reduce la razón de dominancia mediante un corrimiento $k_s$ cercano (y mayor) a $\keff$}
		\item elimina la necesidad de resolver sistemas lineales
		\item modifica el valor de $\keff$ del problema
		\item solo es aplicable a matrices no simétricas
	\end{enumerate}
	
	\item La condición $\varphi(a) = 0$ en una frontera es de tipo:
	\begin{enumerate}[label=\alph*)]
		\item Neumann 
		\item Robin 
		\item periódica 
		\item {\bfseries\boldmath\color{red} Dirichlet}
	\end{enumerate}
	
	\item La frontera reflejante (corriente neta nula, $-D\varphi' = 0$) corresponde a una condición:
	\begin{enumerate}[label=\alph*)]
		\item de Dirichlet homogénea 
		\item {\bfseries\boldmath\color{red} de Neumann homogénea}
		\item de Robin con $\gamma\to\infty$ 
		\item ninguna de las anteriores
	\end{enumerate}
	
	\item Una condición de Robin:
	\begin{enumerate}[label=\alph*)]
		\item fija solo el valor de $\varphi$ 
		\item fija solo el valor de $\varphi'$
		\item {\bfseries\boldmath\color{red} impone una combinación lineal $\alpha\varphi + \beta\varphi' = g$} 
		\item impone periodicidad
	\end{enumerate}
	
	\item La factorización de Cholesky requiere que la matriz sea:
	\begin{enumerate}[label=\alph*)]
		\item {\bfseries\boldmath\color{red} simétrica y definida positiva} 
		\item cualquier matriz no singular
		\item triangular 
		\item diagonal
	\end{enumerate}
	
	\item El costo de la factorización triple de una matriz tridiagonal $n\times n$ es:
	\begin{enumerate}[label=\alph*)]
		\item $\Order{n^3}$ 
		\item $\Order{n^2}$ 
		\item $\Order{n\log n}$ 
		\item {\bfseries\boldmath\color{red} $\Order{n}$}
	\end{enumerate}
	
	\item Para $\lambda h \gg 1$, el factor de amplificación de CN tiende a:
	\begin{enumerate}[label=\alph*)]
		\item $0$ 
		\item {\bfseries\boldmath\color{red} $-1$} 
		\item $+1$ 
		\item $\infty$
	\end{enumerate}
	
	\item Euler explícito para $\vect{\varphi}' = -\mat{A}\vect{\varphi}$ ($\mat{A}$ simétrica definida positiva) es estable si:
	\begin{enumerate}[label=\alph*)]
		\item $\Delta t < 1/\rho(\mat{A})$ 
		\item $\Delta t < \rho(\mat{A})$ 
		\item {\bfseries\boldmath\color{red} $\Delta t < 2/\rho(\mat{A})$} 
		\item $\Delta t < 4/\rho(\mat{A})$
	\end{enumerate}
	
	\item El teorema de Lax--Richtmyer establece que para problemas lineales bien planteados:
	\begin{enumerate}[label=\alph*)]
		\item {\bfseries\boldmath\color{red} consistencia y estabilidad implican convergencia} 
		\item estabilidad implica consistencia
		\item consistencia implica estabilidad 
		\item la convergencia no depende de la estabilidad
	\end{enumerate}
	
	\item Con nodos fantasma y diferencias centrales, la condición $\varphi'(0) = 0$ conduce a:
	\begin{enumerate}[label=\alph*)]
		\item $\varphi_{-1} = 0$ 
		\item $\varphi_{-1} = -\varphi_1$ 
		\item $\varphi_{-1} = \varphi_0$ 
		\item {\bfseries\boldmath\color{red} $\varphi_{-1} = \varphi_1$}
	\end{enumerate}
	
	\item En una malla 2D con $n\times n$ nodos interiores y numeración lexicográfica, el semiancho de banda es:
	\begin{enumerate}[label=\alph*)]
		\item $1$ 
		\item {\bfseries\boldmath\color{red} $n$} 
		\item $n^2$ 
		\item $2$
	\end{enumerate}
	
	\item La iteración $\vect{x}^{(k+1)} = \mat{M}\vect{x}^{(k)} + \vect{c}$ converge para todo $\vect{x}^{(0)}$ si y solo si:
	\begin{enumerate}[label=\alph*)]
		\item $\det\mat{M} < 1$ 
		\item $\rho(\mat{M}) > 1$ 
		\item {\bfseries\boldmath\color{red} $\rho(\mat{M}) < 1$} \item $\mat{M}$ es simétrica
	\end{enumerate}
	
	\item Si la razón de dominancia es cercana a 1:
	\begin{enumerate}[label=\alph*)]
		\item {\bfseries\boldmath\color{red} la convergencia del método de las potencias es lenta} 
		\item la convergencia es más rápida
		\item el método diverge 
		
		\item no tiene efecto
	\end{enumerate}
	
\end{enumerate}